\par\Needspace{8\baselineskip}\bigskip
\pdfbookmark[3]{Exposition, 13}{pallas-II-article13}
\begin{center}13.\end{center}
\noindent\hspace*{\parindent}Suivant l’usage ordinaire, l’époque de longitude moyenne pour le tems \(0\), est \(L-nt\); ainsi suivant cet usage on aurait
\[
\begin{aligned}
&\mathrm{d}(\text{Époque long. m.})=-t\mathrm{d}n+(1-\cos i)\mathrm{d}\AscendingNode\\
&\quad-(2ar+aa\cos\varphi\mathbin{.}\operatorname{tang}\tfrac12\varphi\mathbin{.}\cos(v-\tilde\omega))Tn\mathrm{d}t\\
&\quad+\dfrac{ar\operatorname{tang}\tfrac12\varphi\mathbin{.}\sin(v-\tilde\omega)\mathbin{.}(2+e\cos(v-\tilde\omega))}{\cos\varphi}Vn\mathrm{d}t,
\end{aligned}
\]
où l’on pourrait substituer pour\raisebox{-.45ex}{-} \(\mathrm{d}n\) sa valeur trouvée ci-dessus. Mais de cette manière la variation instantanée de l’époque impliquerait un terme, qui, à cause du facteur \(t\), serait susceptible de croître au delà de toute limite. Pour éviter cet inconvénient, nous définirons l’époque de longitude moyenne par \(L-\int n\mathrm{d}t\), ce qui, comme on voit, revient à la manière ordinaire pour le mouvement elliptique. Nommant donc \(\epsilon\) l’époque de longitude moyenne pour le tems zéro, nous aurons
\[
\begin{array}{rl}
70)\quad&\begin{aligned}[t]\mathrm{d}\epsilon&=(1-\cos i)\mathrm{d}\AscendingNode\\
&\quad-(2ar+aa\cos\varphi\mathbin{.}\operatorname{tang}\tfrac12\varphi\mathbin{.}\cos(v-\tilde\omega))Tn\mathrm{d}t\\
&\quad+\dfrac{ar\operatorname{tang}\tfrac12\varphi\mathbin{.}\sin(v-\tilde\omega)\mathbin{.}(2+e\cos(v-\tilde\omega))}{\cos\varphi}Vn\mathrm{d}t,\end{aligned}
\end{array}
\]
et la longitude moyenne pour le tems \(t\) sera déterminée par la formule
\[
\begin{array}{rl}
71)\quad&L=\epsilon+\int n\mathrm{d}t.
\end{array}
\]
